\section{Three connected research questions}

\textbf{Background.} Cooperation survives on little
information~\cite{morozov2026,leung2026}, but both results assume the
signal is authentic. Authorised push payment fraud breaks that. UK Finance
reports \pounds576.4 million lost across 248{,}070 cases in
2025~\cite{ukfinance2026}, and since 2024 the Payment Systems Regulator
splits reimbursement between sending and receiving
providers~\cite{psr2024}. The bank's only route to intent is the customer,
whose answers the fraudster has scripted: she is a sincere reporter of a
false belief, so a paid check returns text the adversary wrote.

\textbf{Stakeholders.} The payer reads friction as obstruction; the
sending bank carries the complaints; the receiving bank holds the
informative record; the regulator wants fewer losses without
over-blocking. No actor's best choice is defined without the others', so
protection stays low because no institution gains by moving first.

\textbf{Nobel-recognised foundation.} Mechanism design---Hurwicz, Maskin
and Myerson, 2007 Prize in Economic Sciences---asks what an institution
can implement when agents report private
information~\cite{myerson1981}. It assumes a reporter who knows her type
and may misreport. Mine is sincere, does not know her type, and is
scripted by the adversary being screened for.

\textbf{Q1---Economics:} how should a verification budget be allocated
when checking is costly, may be uninformative, and loses effectiveness
with use? \textbf{Q2---Computation:} which mechanisms survive an adversary
who scripts the reporting party? \textbf{Q3---Behavior:} is decay or
authorship the binding constraint?

\textbf{Integrated question.} The three describe one object: a budget that
is scarce, depletes, and must be recomputed against a best-responding
adversary. Alone each reverses the others---concentrating checks on
high-value cases~\cite{benporath2014} is what a script targets, and an
interface remedy for habituation~\cite{anderson2016} cannot reach an
authored signal. Figure~\ref{fig:teaser} states what each discipline
establishes.

\section{Answer to Q1: incentives and intellectual lineage}

\textbf{Hypothesis.} A mechanism is \emph{coaching-proof} if its decision
is invariant to everything the reporting party says within the episode.
H1, falsifiable: no mechanism conditioning on within-episode testimony is
coaching-proof, since a passing script always exists.

\textbf{The game.} The fraudster chooses a script rate $s\in[0,1]$; the
bank chooses a mechanism $m\in\{\text{pay},\text{testimony},
\text{record},\text{hold}\}$ per episode---a mechanism, not a verdict,
since the signal decides; the liability split is inherited from PSR
practice. Information is incomplete in Harsanyi's
sense~\cite{harsanyi1968}: the bank sees only the signal, the payer not
even her own type. Bank payoffs are $-a$ for fraud paid,
$-0.3a$ for a legitimate payment held, $+0.01a$ for one made, and 5 per
record query. The solution concept is Bayesian Nash equilibrium.

\textbf{Benchmark.} Ben-Porath, Dekel and Lipman~\cite{benporath2014}
show that under costly but \emph{reliable} verification the optimal
mechanism concentrates checking, which motivates indexing friction to a
subset of counterparties. Unresolved: checks that are costly \emph{and}
uninformative. Loss figures are observed, coefficients stipulated, the
concentration result inherited.

\section{Answer to Q2: method and evaluation}

\textbf{Inputs and algorithm.} Twelve episodes per round are drawn from a
seeded sequence, each fraudulent with probability 0.25 and, if so,
scripted with probability $s$. Testimony consumes one unit of a
credibility stock of 12 that never refills, and an uncoached payer's
chance of noticing falls with the stock. The record check flags a
fraudulent payee with probability $p$, the cross-institution reporting
rate, and a legitimate one with probability 0.05. Between rounds the
fraudster raises $s$ in proportion to the testimony checks faced.

\textbf{Baseline and modification.} The baseline sweeps $p$ from 0 to 1
across 200 seeds; the metric is the pooled miss rate, fraud paid over
fraud present. The modification switches habituation off and sweeps the
script rate.

\textbf{Obtained.} The miss rate tracks $1-p$ to within 0.032, so
effectiveness is carried by participation, not software. Removing
habituation lowers it from 0.638 to 0.237 at script rate 0, but at script
rate 1 both reach 1.000: habituation is a degradation an interface could
address, authorship a ceiling it cannot.

\textbf{Check.} The notebook re-implements the game's pseudorandom
generator bit-for-bit and reproduces three payoffs recorded by hand in the
deployed game---$-6{,}775$, $-4{,}480$, $-2{,}005$---exactly, so the two
artifacts are one model.

\section{Answer to Q3: behavior and competing explanations}

\textbf{Two explanations.} Anderson et al.~\cite{anderson2016} give
physiological evidence that warnings habituate and propose an interface
remedy. The competitor is authorship: the adversary wrote the answer, so
it never conflicts however much attention it receives.

\textbf{Discriminating evidence.} The two agree at low script rates and
separate at high ones. The ablation is that test in simulation: the gap
between curves is 0.401 at script rate 0 and 0.000 at script rate 1. The
next test runs on people, varying warning novelty and script quality
independently.

\textbf{Limits.} If authorship binds, the bank buys nothing rather than
noisy truth, so optimal allocation places no weight on testimony.
Simulated results establish how the model behaves, not how people do. A
testimony-based mechanism invariant to a passing script---reading latency
rather than content---would falsify H1.

\section{Advanced development and future directions}

\textbf{Foundations and boundary.} The enduring question is trust under
adaptation: keeping exchange credible when those who exploit it adapt.
Mechanism design (Nobel 2007) establishes what an institution can
implement given what agents report~\cite{myerson1981}; interactive proof
systems (Turing 2012)~\cite{goldwasser1989} are the same structure one
level down, a verifier interrogating a possibly adversarial prover. What
is distinctive now is that scripts are cheap and adaptive, so the
reporting channel is adversarially controlled as routine. Current work
treats this as detection; the boundary advanced here is design.

\textbf{Contributions and governed risk.} Economics gains a
costly-verification problem whose cost is a depleting stock, not a
per-unit price; computation gains a decidable property and a recomputation
requirement; behavioral science enters as mechanism state rather than
interface design. Together they explain why mandated warnings keep
expanding after being shown ineffective: their cost falls on a stock
nobody accounts for. The risk to govern is guilt by association on records
payees cannot contest; appeal, expiry and published group false-flag rates
are conditions of the design.

\begin{table}[h]
\caption{From laureate foundations to a new interdisciplinary contribution.}
\label{tab:roadmap}
\scriptsize
\begin{tabularx}{\columnwidth}{@{}lY@{}}
\toprule
Horizon & Milestone and evidence \\
\midrule
Foundation & Mechanism design (Nobel 2007); interactive proofs (Turing 2012). \\
PS1 & Seeded sweep and ablation; notebook reproduces the game's payoffs. \\
Next study & Behavioral test separating habituation from authorship. \\
Longer term & Payee-reporting protocol with appeal, expiry, false-flag rates. \\
2056 & Coaching-proofness adopted as a verification standard. \\
\bottomrule
\end{tabularx}
\end{table}

\noindent\textbf{Expected 2056 Nobel Prize and Turing Award contribution
statement.} By 2056, I aspire to be recognised for a testable criterion of
coaching-proofness and a cross-institution protocol addressing trust under
adaptation in an era when the reporting channel is routinely adversarially
controlled. Building on mechanism design (Hurwicz, Maskin and Myerson,
2007) and interactive proofs (Goldwasser and Micali, 2012), I will
integrate an economics of depleting verification budgets, a computation of
mechanisms recomputed against adaptive adversaries, and a behavioral
account of habituation as mechanism state, to advance the boundary between
detecting deception and designing against it. Progress will be shown by
the milestones in Table~\ref{tab:roadmap}, with intellectual merits in
what verification can conclude from an authored signal and practical
impacts for payers least able to document themselves.

{\small
\subsection*{Open Science Statement}
\textbf{GitHub:} \url{https://github.com/ltishere/PS1-Tian}\quad
\textbf{Colab:} \url{https://colab.research.google.com/github/ltishere/PS1-Tian/blob/main/notebook/coaching_proof_sweep.ipynb}\quad
Commit \texttt{8172a7a}. Dependencies are NumPy and Matplotlib,
preinstalled in Colab; no data is needed. Shortest run: open the Colab
link, Runtime $\rightarrow$ Run all ($\approx$30~s). Stored outputs are
the sweep, the ablation and the check reproducing $-6{,}775$, $-4{,}480$
and $-2{,}005$. Reuse: MIT License.

\subsection*{Statement of Contribution to the UN's SDGs}
\textbf{Hugging Face:}
\url{https://dku-comsci-econ206-2026-tian-liang.static.hf.space/index.html}\quad
This open-access browser game supports \textbf{SDG 4---Quality
Education}~\cite{un2015sdg4} by helping undergraduates learn why a
mechanism's effectiveness depends on who authored the signal it reads.
The learner plays the bank, chooses a mechanism rather than a verdict, varies
the reporting rate, and sees testimony checks drain a credibility stock
while record checks do not. It needs only a browser, and a fixed seed lets
an instructor set identical payments for a class. A feasible evaluation
asks students to predict which mechanism survives a coached payer before
and after play; the indicator is the change in correct
\emph{justifications}. Whether it improves learning remains to be
evaluated.
\par}
\balance